\documentclass[11pt]{article}

\usepackage[letterpaper, top=1in, bottom=1in, left=1.1in, right=1.1in]{geometry}
\usepackage{xcolor}
\usepackage{enumitem}
\usepackage{titlesec}
\usepackage{tcolorbox}
\usepackage{hyperref}
\usepackage{listings}
\usepackage{parskip}

\definecolor{accent}{HTML}{1F4E5F}
\definecolor{lightbg}{HTML}{EAF1F4}
\definecolor{codebg}{HTML}{F2F2F2}

\titleformat{\section}{\Large\bfseries\color{accent}}{\thesection}{0.6em}{}
\titleformat{\subsection}{\large\bfseries\color{accent!80!black}}{\thesubsection}{0.6em}{}
\titlespacing*{\section}{0pt}{1.4em}{0.6em}
\titlespacing*{\subsection}{0pt}{1.1em}{0.4em}

\hypersetup{
	colorlinks=true,
	urlcolor=accent,
	linkcolor=accent,
}

\lstdefinestyle{code}{
	backgroundcolor=\color{codebg},
	basicstyle=\ttfamily\small,
	breaklines=true,
	frame=none,
	xleftmargin=2.2em,
	xrightmargin=0.6em,
	framexleftmargin=2.2em,
	aboveskip=0.6em,
	belowskip=0.8em,
}
\lstset{style=code}

\newtcolorbox{noteBox}{
	colback=lightbg,
	colframe=accent,
	boxrule=0.4pt,
	leftrule=2pt,
	toprule=0pt,
	bottomrule=0pt,
	rightrule=0pt,
	arc=0pt,
	outer arc=0pt,
	boxsep=4pt,
	left=8pt, right=8pt, top=4pt, bottom=4pt,
}

\setlist[enumerate]{itemsep=2pt, topsep=4pt}
\setlist[itemize]{itemsep=2pt, topsep=4pt}

\begin{document}
	
	% ============================================================
	% TITLE PAGE
	% ============================================================
	\begin{titlepage}
		\centering
		\vspace*{2.2in}
		
		{\Huge\bfseries\color{accent} Computer Setup Guide}\\[0.6em]
		{\Large\color{gray!70!black} for Data Analytics Course (Undergraduate)}\\[0.2em]
		{\large\color{gray!60!black} R, RStudio, and course materials}
		
		\vspace{2.5em}
		
		{\large Vojt\v{e}ch Z\'{i}ka, Ph.D.}\\[0.4em]
		{\large \today}
		
		\vfill
	\end{titlepage}
	
	\pagenumbering{gobble}
	\tableofcontents
	\newpage
	\pagenumbering{arabic}
	
	% ============================================================
	% INTRO
	% ============================================================
	This guide walks you through installing everything you need for this course: Git, R, and RStudio. It covers both Windows and Mac --- follow only the steps for your operating system. By the end, you will have the course materials downloaded to your computer and be ready to run R code.
	
	You do not need a GitHub account for this course. The course repository is public, so you can download it directly --- no sign-up required.
	
	\section*{What you will install}
	\begin{itemize}
		\item Git --- the tool used to download (and later update) the course materials.
		\item R --- the programming language used in this course.
		\item RStudio --- the application you will actually work in; it uses R behind the scenes.
	\end{itemize}
	Install them in this order: Git, then R, then RStudio. This takes about 15--20 minutes total.
	

	\section{Install Git}
	
	\subsection*{Windows}
	\begin{enumerate}
		\item Go to \url{https://git-scm.com/download/win} --- the download should start automatically.
		\item Run the installer. You can accept all the default options by repeatedly clicking ``Next,'' then ``Install.''
		\item Once finished, open the Start menu, search for ``Git Bash,'' and open it.
		\item Type the following and press Enter:
\begin{lstlisting}
git --version
\end{lstlisting}
	\end{enumerate}
	If you see a version number (e.g. \texttt{git version 2.44.0}), Git is installed correctly.
	
	\subsection*{Mac}
	\begin{enumerate}
		\item Open the Terminal app (press Cmd+Space, type ``Terminal,'' press Enter).
		\item Type the following and press Enter:
\begin{lstlisting}
git --version
\end{lstlisting}
		\item If Git is not installed yet, macOS will show a popup asking to install the ``command line developer tools.'' Click Install and wait a few minutes for it to finish.
		\item Once it is done, run \texttt{git --version} again to confirm you see a version number.
	\end{enumerate}
	
	\begin{noteBox}
		\textbf{Note:} If no popup appears and you get a ``command not found'' error instead, download Git directly from \url{https://git-scm.com/download/mac} and run the installer.
	\end{noteBox}
	

	\section{Install R}
	\begin{enumerate}
		\item Go to \url{https://cran.r-project.org}.
		\item Windows: click ``Download R for Windows'' $\rightarrow$ ``base'' $\rightarrow$ the top download link (e.g. ``Download R-4.x.x for Windows'').
		\item Mac: click ``Download R for macOS'' and choose the correct file for your Mac --- the page lists an ``arm64'' build (for Apple Silicon Macs: M1/M2/M3/M4) and an ``x86\_64'' build (for older Intel Macs). If you are not sure which chip your Mac has, check Apple menu $\rightarrow$ About This Mac.
		\item Run the installer and accept all default options.
	\end{enumerate}
	
	\begin{noteBox}
		\textbf{Note:} You will not open R directly in this course --- you will always work through RStudio, which uses R automatically in the background.
	\end{noteBox}
	

	\section{Install RStudio}
	\begin{enumerate}
		\item Go to \url{https://posit.co/download/rstudio-desktop/}.
		\item Scroll to ``Install RStudio Desktop'' and download the free version for your operating system.
		\item Run the installer and accept all default options.
		\item Open RStudio once to confirm it launches correctly. You should see a window with several panels (Console, Environment, Files, etc.).
	\end{enumerate}
	

\section{Download the course materials}
The course repository is public, so you can download it without creating any account. You will use Git to ``clone'' it --- this downloads the whole folder structure (scripts, data, slides) to your computer.
\begin{enumerate}
	\item Choose (or create) a folder on your computer where you want the course materials to live. For example, create a new folder called \texttt{DataAnalytics} on your Desktop. (Avoid spaces in the folder name --- they cause problems in the next step.)
	\item Get the full path to that folder:
	\begin{itemize}
		\item \textbf{Mac:} in Finder, hold the \texttt{Option} key and right-click (or Control-click) the folder, then choose ``Copy `\texttt{DataAnalytics}' as Pathname.''
		\item \textbf{Windows:} in File Explorer, hold \texttt{Shift} and right-click the folder, then choose ``Copy as path.''
	\end{itemize}
	\item Open Terminal (Mac) or Git Bash (Windows). Type \texttt{cd } (with a space after it), then paste the path you just copied, and press Enter. For example, on Mac:
\begin{lstlisting}
cd /Users/yourname/Desktop/DataAnalytics
\end{lstlisting}
	On Windows, ``Copy as path'' gives a \texttt{C:\textbackslash{}...} style path wrapped in quotes --- paste it as-is, quotes included:
\begin{lstlisting}
cd "C:\Users\yourname\Desktop\DataAnalytics"
\end{lstlisting}
	\item Clone the repository:
\begin{lstlisting}
git clone https://github.com/vojtechzika/data-analytics-ug.git
\end{lstlisting}
	\item This creates a new folder called \texttt{data-analytics-ug} inside the folder you selected, containing all the course files.
\end{enumerate}

	\section{Open the project in RStudio}
	\begin{enumerate}
		\item Open the \texttt{data-analytics-ug} folder you just downloaded, in Finder (Mac) or File Explorer (Windows).
		\item Double-click \texttt{data-analytics-ug.Rproj}.
		\item RStudio should open with this project loaded --- you will see ``data-analytics-ug'' in the top-right corner of the RStudio window.
	\end{enumerate}

	\section{Verify everything works}
	\begin{enumerate}
		\item In RStudio, click into the Console panel (usually bottom-left).
		\item Type the following and press Enter:
\begin{lstlisting}
1 + 1
\end{lstlisting}
		\item You should see \texttt{[1] 2} printed below.
		\item In the Files panel (bottom-right), confirm you can see the \texttt{R}, \texttt{data}, and \texttt{slides} folders.
	\end{enumerate}
	If all of the above worked, your setup is complete.
	

	\section{Getting course updates later}
	Throughout the course, new materials may be added to the repository. To get the latest version, open Terminal or Git Bash, navigate into your course folder, and run:
\begin{lstlisting}
cd Documents/data-analytics-ug
git pull
\end{lstlisting}
	This downloads any new or updated files without affecting your own work, as long as you have not edited the same files the instructor changed.
	
	\newpage
	\section{Troubleshooting}
	\begin{itemize}
		\item \textbf{``git is not recognized'' / ``command not found''} right after installing: close and reopen Terminal or Git Bash (or restart your computer), then try again.
		\item \textbf{Nothing happens when double-clicking the \texttt{.Rproj} file:} make sure RStudio is installed, then right-click the file and choose ``Open with $\rightarrow$ RStudio.''
		\item \textbf{\texttt{git clone} fails with a permissions or authentication error:} double-check the URL was typed exactly as shown --- since the repo is public, no login should be requested.
	\end{itemize}
	
\end{document}